pdfoutput=1
% Options for packages loaded elsewhere
\PassOptionsToPackage{unicode}{hyperref}
\PassOptionsToPackage{hyphens}{url}
\pdfoutput=1
\documentclass[
  12pt,
]{article}
\usepackage{xcolor}
\usepackage{amsmath,amssymb}
\setcounter{secnumdepth}{-\maxdimen} % remove section numbering
\usepackage{iftex}
\ifPDFTeX
  \usepackage[T1]{fontenc}
  \usepackage{textcomp} % provide euro and other symbols
\else % if luatex or xetex
  \usepackage{unicode-math} % this also loads fontspec
  \defaultfontfeatures{Scale=MatchLowercase}
  \defaultfontfeatures[\rmfamily]{Ligatures=TeX,Scale=1}
\fi
\usepackage{lmodern}
\ifPDFTeX\else
  % xetex/luatex font selection
\fi
% Use upquote if available, for straight quotes in verbatim environments
\IfFileExists{upquote.sty}{\usepackage{upquote}}{}
\IfFileExists{microtype.sty}{% use microtype if available
  \usepackage[]{microtype}
  \UseMicrotypeSet[protrusion]{basicmath} % disable protrusion for tt fonts
}{}
\makeatletter
\@ifundefined{KOMAClassName}{% if non-KOMA class
  \IfFileExists{parskip.sty}{%
    \usepackage{parskip}
  }{% else
    \setlength{\parindent}{0pt}
    \setlength{\parskip}{6pt plus 2pt minus 1pt}}
}{% if KOMA class
  \KOMAoptions{parskip=half}}
\makeatother
\usepackage{longtable,booktabs,array}
\usepackage{calc} % for calculating minipage widths
% Correct order of tables after \paragraph or \subparagraph
\usepackage{etoolbox}
\makeatletter
\patchcmd\longtable{\par}{\if@noskipsec\mbox{}\fi\par}{}{}
\makeatother
% Allow footnotes in longtable head/foot
\IfFileExists{footnotehyper.sty}{\usepackage{footnotehyper}}{\usepackage{footnote}}
\makesavenoteenv{longtable}
\setlength{\emergencystretch}{3em} % prevent overfull lines
\providecommand{\tightlist}{%
  \setlength{\itemsep}{0pt}\setlength{\parskip}{0pt}}
\usepackage{bookmark}
\IfFileExists{xurl.sty}{\usepackage{xurl}}{} % add URL line breaks if available
\urlstyle{same}
\hypersetup{
  hidelinks,
  pdfcreator={LaTeX via pandoc}}

\author{SCX}
\date{2026}

\begin{document}

\section{Fixed Point and Convergence Theorem of SCX
Self-Evolution}\label{fixed-point-and-convergence-theorem-of-scx-self-evolution}

\begin{quote}
\textbf{Version}: 2026-06-28 \textbar{} \textbf{Status}: Theoretical
framework (partial proof) \textbar{} \textbf{Scope}: Central theorem
establishing convergence regimes for the self-evolution loop
\textbf{Notation note}: \(\mathcal{S}\) denotes the state space
(existing SCX theory). \(S_t\) denotes the gatekeeper scoring function
at time \(t\). \(\mathcal{M}_t\) denotes the memory bank.
\(f_{\theta_t}\) denotes the NEP student.
\end{quote}

\begin{center}\rule{0.5\linewidth}{0.5pt}\end{center}

\subsection{Table of Contents}\label{table-of-contents}

\begin{enumerate}
\def\labelenumi{\arabic{enumi}.}
\tightlist
\item
  \hyperref[1-introduction-the-central-question]{Introduction: The
  Central Question}
\item
  \hyperref[2-theorem-se-1-convergence-of-scx-self-evolution-formal-statement]{Theorem
  SE-1: Convergence of SCX Self-Evolution (Formal Statement)}
\item
  \hyperref[3-conditions-for-theorem-se-1]{Conditions for Theorem SE-1}
\item
  \hyperref[4-conclusion-of-theorem-se-1]{Conclusion of Theorem SE-1}
\item
  \hyperref[5-lemma-se-11-lyapunov-function-non-increase]{Lemma SE-1.1:
  Lyapunov Function Non-Increase}
\item
  \hyperref[6-lemma-se-12-finite-cover-properties-of-memory-bank-types]{Lemma
  SE-1.2: Finite Cover Properties of Memory Bank Types}
\item
  \hyperref[7-lemma-se-13-parameter-displacement-vanishes]{Lemma SE-1.3:
  Parameter Displacement Vanishes}
\item
  \hyperref[8-lemma-se-14-limit-points-are-fixed-points]{Lemma SE-1.4:
  Limit Points Are Fixed Points}
\item
  \hyperref[9-proof-sketch-of-theorem-se-1]{Proof Sketch of Theorem
  SE-1}
\item
  \hyperref[10-four-convergence-regimes]{Four Convergence Regimes}
\item
  \hyperref[11-conditions-separating-the-regimes]{Conditions Separating
  the Regimes}
\item
  \hyperref[12-connection-to-theorem-1-noise-detection]{Connection to
  Theorem 1 (Noise Detection)}
\item
  \hyperref[13-connection-to-theorem-3-unidentifiability]{Connection to
  Theorem 3 (Unidentifiability)}
\item
  \hyperref[14-what-is-proven-vs-conjectured]{What is Proven
  vs.~Conjectured}
\item
  \hyperref[15-references]{References}
\end{enumerate}

\begin{center}\rule{0.5\linewidth}{0.5pt}\end{center}

\subsection{1. Introduction: The Central
Question}\label{introduction-the-central-question}

The self-evolution loop couples three components:

\[
\begin{aligned}
\mathcal{M}_t &\gets \text{AcceptOrReject}(\mathcal{D}_t, S_t) \quad &\text{(gatekeeping)} \\
\theta_{t+1} &\gets \theta_t - \alpha_t \nabla_\theta \ell(f_{\theta_t}, \mathcal{M}_t) \quad &\text{(student update)} \\
S_{t+1} &\gets \text{UpdateSCX}(S_t, \mathcal{M}_t, f_{\theta_{t+1}}) \quad &\text{(gatekeeper update)}
\end{aligned}
\]

\textbf{Central question}: Does the sequence
\((S_t, \theta_t, \mathcal{M}_t)\) converge? If so, to what? If not,
what alternative long-term behaviors are possible?

This document provides the central theorem (SE-1) that characterizes the
convergence properties of the self-evolution system.

\begin{center}\rule{0.5\linewidth}{0.5pt}\end{center}

\subsection{2. Theorem SE-1: Convergence of SCX Self-Evolution (Formal
Statement)}\label{theorem-se-1-convergence-of-scx-self-evolution-formal-statement}

\textbf{Theorem SE-1 (Convergence of SCX Self-Evolution).} Let
\((S_t, \theta_t, \mathcal{M}_t)\) be the sequence generated by the SCX
self-evolution algorithm:

\[
\begin{aligned}
\mathcal{M}_{t+1} &= \mathcal{M}_t \cup \{(x, y) \in \mathcal{D}_{t+1} : S_t(x,y) \geq \gamma\} \\
\theta_{t+1} &= \theta_t - \alpha_t \nabla_\theta \ell(f_{\theta_t}(x_t), y_t), \quad (x_t, y_t) \sim \text{Unif}(\mathcal{M}_t) \\
S_{t+1} &= \Pi_{[0,1]} \left[ S_t + \beta_t \cdot \bigl( \text{SCXUpdate}(S_t, \mathcal{M}_{t+1}, f_{\theta_{t+1}}) - S_t \bigr) \right]
\end{aligned}
\]

where \(\gamma \in (0,1)\) is the acceptance threshold, \(\alpha_t\) is
the student learning rate, \(\beta_t\) is the gatekeeper update rate,
and \(\Pi_{[0,1]}\) projects scores to \([0,1]\).

\begin{quote}
\textbf{Reference-set evaluation (C10).} The gatekeeper update
\(\text{SCXUpdate}(S_t, \mathcal{M}_{t+1}, f_{\theta_{t+1}})\) is
evaluated on a fixed reference set \(\mathcal{M}_0\) (see condition
C10), not on the growing memory bank. This anchors the evaluation to a
stationary distribution and is required for Lyapunov descent (Lemma
SE-1.1). Without reference-set replay, the gatekeeper can trivially
reduce its loss by accepting only easy samples (Theorem 12.2).
\end{quote}

Assume the following conditions hold:

\textbf{(C1) Finite structure space.} The input space \(\mathcal{X}\) is
finite or has finite covering dimension: \(\dim(\mathcal{X}) < \infty\).

\textbf{\textgreater{} FIXED (2026-06-28, DEFECT-15):} The original
``finite \(\mathcal{X}\)'' assumption is violated in practice
(continuous \(\mathbb{R}^{d_\phi}\) feature spaces). This section
replaces it with a \textbf{covering number argument} that achieves the
same memory bank stabilization result (Lemma SE-1.2) under the much
weaker condition of finite covering dimension. See Section 3.1 for the
full construction.

\textbf{(C2) Lipschitz student.} The NEP student \(f_\theta\) is
Lipschitz continuous in \(\theta\) with constant \(L_f\):

\[
\|f_{\theta_1}(x) - f_{\theta_2}(x)\| \leq L_f \|\theta_1 - \theta_2\|, \quad \forall x \in \mathcal{X}, \forall \theta_1, \theta_2
\]

\textbf{(C3) Lipschitz gatekeeper.} The SCX update function is Lipschitz
continuous with constant \(L_S\):

\[
\|\text{SCXUpdate}(S_1, \mathcal{M}, f) - \text{SCXUpdate}(S_2, \mathcal{M}, f)\|_\infty \leq L_S \|S_1 - S_2\|_\infty
\]

\textbf{(C4) Learning rates.} The student learning rate \(\alpha_t\)
satisfies the Robbins-Monro conditions:

\[
\sum_{t=1}^\infty \alpha_t = \infty, \quad \sum_{t=1}^\infty \alpha_t^2 < \infty
\]

The gatekeeper update rate \(\beta_t\) satisfies:

\[
\beta_t \to 0, \quad \sum_{t=1}^\infty \beta_t^2 < \infty
\]

\begin{quote}
\textbf{FIXED (2026-06-28, B8):} The original C4 listed
\(\sum \beta_t = \infty\) (Robbins-Monro exploration) for the
gatekeeper. However, for the cumulative distribution shift to be finite
(needed for Lyapunov descent), the stronger condition
\(\sum_{t=1}^\infty \beta_t < \infty\) is required. This is now stated
as an \textbf{explicit separate sub-condition (C4b)} --- it does NOT
follow from \(\beta_t = o(\alpha_t)\) (C6) combined with
\(\sum \alpha_t^2 < \infty\). See the counterexample in the C6
discussion below.

\textbf{Resolution of the tension:} The gatekeeper exploration
(originally \(\sum \beta_t = \infty\)) is instead provided by the random
exploration fraction (C9) and the annealing acceptance threshold (C8).
The gatekeeper update rate itself must satisfy \(\sum \beta_t < \infty\)
to ensure finite total distribution shift. The canonical schedule is
updated to \(\alpha_t = t^{-a}\), \(\beta_t = t^{-b}\) with
\(\frac{1}{2} < a < 1 < b\), e.g., \(a = 0.6\), \(b = 1.2\).
\end{quote}

\textbf{(C5) Conditional i.i.d. sampling.} Conditioned on \(S_t\), the
samples drawn at time \(t\) are i.i.d. from the acceptance-biased
distribution:

\[
(x_t, y_t) \mid S_t \sim P_{S_t}(x,y) \propto S_t(x,y) \cdot P_0(x,y)
\]

\textbf{(C6) Two-timescale separation (corrected --- replaces original
``sufficient annealing'').} The gatekeeper update rate \(\beta_t\) must
be asymptotically slower than the student learning rate \(\alpha_t\):

\[
\boxed{\;\beta_t = o(\alpha_t)\;}, \qquad \text{i.e.,} \quad \lim_{t\to\infty} \frac{\beta_t}{\alpha_t} = 0.
\]

\textbf{Rationale (DEFECT-14 fix).} The original condition C6
(\(\max(\alpha_t, \beta_t) \to 0\)) is insufficient because it does not
resolve the tension between gatekeeper exploration
(\(\sum \beta_t = \infty\)) and student stabilization (requires
approximately stationary distribution). The two-timescale condition
\(\beta_t = o(\alpha_t)\) resolves this: between any two gatekeeper
updates, the student takes asymptotically many gradient steps, ensuring
local convergence to the current distribution's minimum before the
distribution shifts.

\textbf{Jointly satisfiable scheduling scheme (corrected B8).} The
conditions \(\sum \alpha_t = \infty\), \(\sum \alpha_t^2 < \infty\),
\(\sum \beta_t < \infty\), \(\sum \beta_t^2 < \infty\), and
\(\beta_t = o(\alpha_t)\) are jointly satisfiable. A canonical example
is:

\[\alpha_t = t^{-a}, \quad \beta_t = t^{-b}, \quad \text{with} \quad \frac{1}{2} < a < 1 < b.\]

For \(a = 0.6\), \(b = 1.2\): - \(\sum \alpha_t = \infty\) (harmonic
series with exponent \(< 1\)) -
\(\sum \alpha_t^2 = \sum t^{-1.2} < \infty\) (Robbins-Monro for student)
- \(\sum \beta_t = \sum t^{-1.2} < \infty\) (finite total distribution
shift --- explicit condition) -
\(\sum \beta_t^2 = \sum t^{-2.4} < \infty\) (variance control for
gatekeeper) - \(\beta_t/\alpha_t = t^{-0.6} \to 0\) (two-timescale
separation)

\textbf{Consequence of violation.} Without two-timescale separation
(\(\beta_t\) not \(o(\alpha_t)\)), the cumulative distribution shift may
diverge:

\[\sum_{t=1}^\infty TV(P_{t+1}, P_t) \leq \sum_{t=1}^\infty \frac{2\beta_t B_S}{\mathbb{E}_{P_0}[S_t] - \beta_t B_S}.\]

\textbf{FIXED (2026-06-28, B8): \(\sum \beta_t < \infty\) is a separate
condition, NOT implied by \(\beta_t = o(\alpha_t)\).} The original text
claimed that \(\beta_t = o(\alpha_t)\) combined with
\(\sum \alpha_t^2 < \infty\) implies \(\sum \beta_t < \infty\). This is
\textbf{mathematically false}, as shown by the following counterexample:

\[\alpha_t = t^{-0.6}, \quad \beta_t = \frac{1}{t^{0.6} \log t}.\]

Then: - \(\beta_t / \alpha_t = 1/\log t \to 0\), so
\(\beta_t = o(\alpha_t)\) ✓ -
\(\sum \alpha_t^2 = \sum t^{-1.2} < \infty\) ✓ - But
\(\sum \beta_t = \sum \frac{1}{t^{0.6} \log t} = \infty\) (diverges by
the integral test) ✗

The flaw: \(\beta_t = o(\alpha_t)\) only guarantees
\(\beta_t < \alpha_t\) for large \(t\), but \(\sum \alpha_t^2 < \infty\)
controls the \textbf{squares} of \(\alpha_t\), not \(\alpha_t\) itself.
Since \(\sum \alpha_t = \infty\) (Robbins-Monro), \(\alpha_t\) decays as
\(t^{-a}\) with \(a \in (1/2, 1]\), and \(\beta_t = o(\alpha_t)\) only
forces \(b > a\), not \(b > 1\) (which is what \(\sum \beta_t < \infty\)
requires).

\textbf{Resolution.} The condition
\(\sum_{t=1}^\infty \beta_t < \infty\) is now stated as an
\textbf{explicit separate requirement} in C4 (below), independent of the
two-timescale condition C6. It must be verified independently or
enforced through the specific schedule \(\beta_t = t^{-b}\) with
\(b > 1\) (which is stronger than the Robbins-Monro
\(\sum \beta_t = \infty\) originally listed in C4 --- see the corrected
C4 below).

Under the corrected C4 (which now requires \(\sum \beta_t < \infty\) for
the gatekeeper, distinct from the student's \(\sum \alpha_t = \infty\)),
the total distribution shift is finite and the student sees an
asymptotically stationary distribution.

\textbf{(C7) Gatekeeper bounded gradient.} The SCX update increment is
bounded:

\[
\|\text{SCXUpdate}(S, \mathcal{M}, f) - S\|_\infty \leq B_S < \infty
\]

\textbf{(C8) Annealing acceptance threshold (sufficient, not
necessary).} The acceptance threshold \(\gamma_t\) starts low
(permissive, \(\gamma_0 < 0.5\)) and increases to \(0.5\) as
\(t \to \infty\):

\[\gamma_t = \gamma_0 + (0.5 - \gamma_0)(1 - e^{-t/\tau_\gamma}).\]

This prevents early over-confidence while allowing increasing
selectivity as the gatekeeper improves.

\textbf{(C9) Random exploration fraction (sufficient, not necessary).} A
fixed fraction \(\varepsilon_{\text{explore}} > 0\) of accepted samples
are chosen uniformly at random regardless of \(S_t\) score, ensuring all
regions remain represented in \(\mathcal{M}_t\).

\begin{quote}
\textbf{B7 FIX --- \(\varepsilon > 0\) lower bound for importance
weights:} The importance weight bound \(\|w\|_\infty \leq W < \infty\)
in Theorem 12.5 requires \(P_{S_t}(x) > 0\) for all \(x\) in the support
of \(P_0\), i.e.,
\(\inf_{x \in \operatorname{supp}(P_0)} P_{S_t}(x) > 0\). When
\(\varepsilon_{\text{explore}} = 0\) (no exploration), the gatekeeper
can zero out \(S_t(x)\) for regions it deems unreliable, making
\(P_{S_t}(x) = 0\) and \(w(x) = P_0(x)/P_{S_t}(x) = \infty\) --- the
importance weights become unbounded.

\textbf{Resolution:} Condition C9 requires
\(\varepsilon_{\text{explore}} > 0\) \textbf{strictly}. With
\(\varepsilon_{\text{explore}} > 0\), every \(x\) with \(P_0(x) > 0\)
has \(P_{S_t}(x) \geq \varepsilon_{\text{explore}} \cdot P_0(x) > 0\)
(the uniform exploration component ensures non-zero probability).
Consequently:
\[\|w\|_\infty = \sup_{x} \frac{P_0(x)}{P_{S_t}(x)} \leq \frac{1}{\varepsilon_{\text{explore}}} = W < \infty.\]
The bound \(W\) is finite and explicit:
\(W = 1/\varepsilon_{\text{explore}}\). As
\(\varepsilon_{\text{explore}} \to 0\), \(W \to \infty\), and the
variance of the importance-weighted gradient estimator diverges ---
confirming that \textbf{\(\varepsilon_{\text{explore}} > 0\) is a
necessary condition} for the Lyapunov descent claim in Theorem 12.5 to
be non-vacuous.
\end{quote}

\begin{quote}
\textbf{Note on C8-C9}: These are sufficient conditions to mitigate the
selection bias cycle (§5.1). They are not mathematically necessary for
the theorem statement, but without them (or equivalent mechanisms), the
Lyapunov descent claim in Lemma SE-1.1 is not guaranteed --- the
decrease in \(V_t\) on \(P_{S_t}\) may reflect distribution shift rather
than genuine improvement.
\end{quote}

\textbf{(C10) Reference-set replay (REQUIRED for Lyapunov descent).} A
fixed reference set \(\mathcal{M}_0\) and validation set
\(\mathcal{V}_0\) are held out before self-evolution begins. The
Lyapunov function \(\Psi\) is evaluated on these fixed sets, NOT on the
growing memory bank \(\mathcal{M}_t\):

\[\Psi(S_t, \theta_t) = \frac{1}{|\mathcal{M}_0|} \sum_{x \in \mathcal{M}_0} (S_t(x, y(x)) - \hat{C}(x))^2 + \frac{\lambda}{|\mathcal{V}_0|} \sum_{(x,y) \in \mathcal{V}_0} \ell(f_{\theta_t}(x), y).\]

\begin{quote}
\textbf{Critical dependency}: Condition C10 is the bridge between Lemma
SE-1.1 and Theorems 12.2/12.5. Theorem 12.2 proves that Lyapunov descent
is \textbf{formally impossible} without reference-set replay (the
gatekeeper can trivially reduce loss by accepting only easy samples).
Theorem 12.5 proves that \textbf{with} reference-set replay and
importance sampling, Lyapunov descent is achievable. Lemma SE-1.1 is
therefore \textbf{conditional on C10} --- its ``CONDITIONAL'' status in
the table below assumes C10 holds. \textbf{Without replay, Lemma SE-1.1
is not guaranteed (see Thm 12.2).}

Without C10, Lemma SE-1.1's supermartingale argument is valid in form,
but the Lyapunov function
\(V_t = \mathbb{E}_{P_{S_t}}[\ell(f_{\theta_t}(x), y)]\) measures
performance on the acceptance-biased distribution, which can decrease
even as true performance degrades (the selection bias cycle, §5.1). C10
breaks this cycle by anchoring evaluation to a fixed reference.
\end{quote}

\begin{center}\rule{0.5\linewidth}{0.5pt}\end{center}

\subsection{3. Conditions for Theorem
SE-1}\label{conditions-for-theorem-se-1}

\subsubsection{3.1 Condition C1: Finite Structure Space --- Covering
Number Replacement (DEFECT-15
Fix)}\label{condition-c1-finite-structure-space-covering-number-replacement-defect-15-fix}

\textbf{Original (defective):} C1 assumed \(\mathcal{X}\) is finite.
This is violated in practice for continuous feature spaces
\(\phi: \mathcal{X} \to \mathbb{R}^{d_\phi}\), where the input space is
a continuum and the memory bank can grow without bound in terms of
distinguishable configurations.

\textbf{FIXED (2026-06-28): Covering Number Argument.} We replace the
finite-\(\mathcal{X}\) assumption with a covering-number construction
that achieves the same stabilization result under the much weaker
condition of finite covering dimension.

\textbf{Construction (ε-net partitioning):}

\begin{enumerate}
\def\labelenumi{\arabic{enumi}.}
\item
  \textbf{Finite covering dimension.} Assume the feature space
  \(\Phi = \phi(\mathcal{X}) \subseteq \mathbb{R}^{d_\phi}\) has finite
  Euclidean dimension \(d_\phi < \infty\) (always true for any practical
  feature map).
\item
  \textbf{Lipschitz gatekeeper (C3).} The SCX score function
  \(S_t(x, y)\) is \(L_S\)-Lipschitz in \(x\) by condition C3. This
  means: within any \(\varepsilon/L_S\)-ball in feature space, \(S_t\)
  varies by at most \(\varepsilon\).
\item
  \textbf{ε-net construction.} Partition \(\Phi\) into \(N_\varepsilon\)
  hypercubes of side \(\varepsilon/L_S\):
  \[N_\varepsilon = \left\lceil\frac{\operatorname{diam}(\Phi) \cdot L_S}{\varepsilon}\right\rceil^{d_\phi} = O\!\left(\left(\frac{L_S}{\varepsilon}\right)^{d_\phi}\right)\]
\item
  \textbf{Finite distinguishable types.} Within each cube, all points
  are \(\varepsilon\)-equivalent in gatekeeper score. The memory bank
  state is entirely determined by \textbf{which cubes have been visited
  and accepted}. Since there are \(N_\varepsilon\) cubes and each is
  either visited-or-not, the number of possible memory bank
  configurations is at most:
  \[|\mathcal{M}\text{-configs}| \leq 2^{N_\varepsilon} = \exp\!\left(O\!\left(\left(\frac{L_S}{\varepsilon}\right)^{d_\phi}\right)\right)\]
  which is \textbf{finite} for any fixed \(\varepsilon > 0\).
\item
  \textbf{Stabilization up to ε-precision.} Lemma SE-1.2's finite-cover
  argument now applies: if two memory bank configurations are
  \(\varepsilon\)-close (same cubes visited), the gatekeeper scores
  \(S_t\) differ by at most \(O(\varepsilon)\). Since machine precision
  is finite (\(\varepsilon_{\text{mach}} \approx 10^{-16}\) for
  float64), taking \(\varepsilon = \varepsilon_{\text{mach}}\) yields a
  genuinely finite configuration space for all practical purposes.
\item
  \textbf{Infinite-space extension (sketch).} For the limit
  \(\varepsilon \to 0\), the configuration space grows as
  \(2^{O((1/\varepsilon)^{d_\phi})}\), which is doubly exponential in
  \(1/\varepsilon\). However, the stabilization result only requires
  that \textbf{for any given precision \(\varepsilon > 0\), the system
  enters and stays within an \(\varepsilon\)-neighborhood of a fixed
  point after finite time}. This follows because:

  \begin{itemize}
  \tightlist
  \item
    The Lyapunov descent (conjectured) drives the system toward the
    fixed point
  \item
    Once within distance \(\varepsilon\), the \(\varepsilon\)-net only
    needs to resolve \(O((1/\varepsilon)^{d_\phi})\) cubes
  \item
    Finite-time convergence to \(\varepsilon\)-precision holds for each
    fixed \(\varepsilon\)
  \end{itemize}
\end{enumerate}

\textbf{Practical interpretation.} This covering-number construction
replaces the unrealistic ``finite \(\mathcal{X}\)'' assumption with the
genuinely weak ``finite covering dimension'' assumption, which holds for
\textbf{any} feature space embedded in \(\mathbb{R}^{d_\phi}\) ---
including continuous pixel spaces, molecular descriptor spaces, and
language model embedding spaces. The same stabilization result (Lemma
SE-1.2) holds, now justified by the Lipschitz property (C3) and the
finite-dimensional feature embedding rather than by an artificial
finiteness constraint.

\textbf{Remaining gap.} The infinite-space limit \(\varepsilon \to 0\)
with full functional convergence (rather than \(\varepsilon\)-precision
approximate convergence) requires additional compactness conditions
(e.g., Arzelà-Ascoli for the space of score functions) that are beyond
the current scope. This is noted as an open problem in the verification
report.

\subsubsection{3.2 Conditions C2-C3: Lipschitz
Continuity}\label{conditions-c2-c3-lipschitz-continuity}

These conditions ensure that small changes in parameters or gatekeeper
scores produce proportionally small changes in the system state.

\paragraph{3.2.1 Lemma SE-1.0: Indicator Discontinuity Is Measure-Zero
(B9
Fix)}\label{lemma-se-1.0-indicator-discontinuity-is-measure-zero-b9-fix}

\textbf{Lemma SE-1.0 (Lipschitz-Almost-Everywhere Consensus).} Under
condition C5 (conditional i.i.d. sampling), the consensus score
\(C(x) = \frac{1}{M}\sum_{m=1}^M \mathbf{1}\{\ell(f_m(x), y) > \tau\}\)
is Lipschitz-almost-everywhere with respect to the acceptance-biased
data distribution \(P_{S_t}\). Specifically, the set of points where
\(C\) is discontinuous has \(P_{S_t}\)-measure zero:

\[\mathbb{P}_{(x,y) \sim P_{S_t}}\bigl( \exists m: \ell(f_m(x), y) = \tau \bigr) = 0.\]

\textbf{Why this resolves the B9 gap.} The ARXIV verification identified
a fundamental gap: condition C3 requires the gatekeeper and its update
to be Lipschitz continuous, but the SCX consensus computation \(C(x)\)
uses indicator functions \(\mathbf{1}\{\ell > \tau\}\) that are
discontinuous at the threshold \(\tau\). However:

\begin{enumerate}
\def\labelenumi{\arabic{enumi}.}
\item
  \textbf{Measure-zero under C5.} Since
  \(P_{S_t} \propto S_t \cdot P_0\) is absolutely continuous with
  respect to \(P_0\), and under the mild assumption that
  \(\ell(f_m(x), y)\) has a continuous distribution (true when the
  student \(f_\theta\) is a smooth function --- e.g., any neural network
  with continuous activations --- and the data distribution has a
  density), the event \(\ell(f_m(x), y) = \tau\) has probability zero.
  The discontinuities occur on a set of Lebesgue measure zero.
\item
  \textbf{Lipschitz-almost-everywhere is sufficient.} The Lyapunov
  analysis in Lemma SE-1.1 uses expectations
  \(\mathbb{E}_{P_{S_t}}[\cdot]\), and expectations are invariant to
  modifications on measure-zero sets. The standard theory of stochastic
  approximation (Borkar, 2008; Kushner \& Yin, 2003) requires Lipschitz
  continuity of the mean field, not pointwise continuity. The mean field
  \(\mathbb{E}_{P_{S_t}}[C(x)]\) is Lipschitz continuous even when
  \(C(x)\) has jump discontinuities on a null set, as long as the
  distribution has a density.
\item
  \textbf{Sigmoid smoothing as optional refinement.} For implementations
  concerned with gradient-based optimization through \(C(x)\), the
  indicator can be replaced by a smooth sigmoid
  \(\sigma(\beta(\ell - \tau))\) with \(\beta \gg 1\). The approximation
  error satisfies
  \(\|\mathbf{1}\{\ell > \tau\} - \sigma(\beta(\ell - \tau))\|_1 = O(1/\beta)\).
  Taking \(\beta \to \infty\) recovers the hard threshold, and the proof
  chain remains valid because (a) the expectation converges uniformly in
  \(\beta\), and (b) the Lipschitz constant of the sigmoid-smoothed
  consensus is \(O(\beta)\), which can be absorbed into the step-size
  conditions via the two-timescale separation (C6).
\end{enumerate}

\textbf{Formal justification (measure-zero argument):}

\emph{Proof of Lemma SE-1.0.} For each expert \(m\), the student loss
\(\ell(f_m(x), y)\) is a random variable induced by the sampling
process. Under C5, \((x, y)\) is drawn from \(P_{S_t}\), which has a
density with respect to the base measure. The function
\(g_m(x, y) = \ell(f_m(x), y)\) is continuous in \((x, y)\) (composition
of the continuous student \(f_{\theta_t}\) and the continuous loss
\(\ell\)). For a continuous random variable, the probability of hitting
any specific value \(\tau\) is zero:
\(\mathbb{P}(g_m(X, Y) = \tau) = 0\). By the union bound over the \(M\)
experts:

\[\mathbb{P}\bigl( \exists m: \ell(f_m(X), Y) = \tau \bigr) \leq \sum_{m=1}^M \mathbb{P}\bigl( \ell(f_m(X), Y) = \tau \bigr) = 0.\]

Therefore, \(C(x)\) is almost surely continuous --- its discontinuity
set (the union of the \(M\) level sets \(\{\ell(f_m(x), y) = \tau\}\))
has probability zero. On the complement of this null set, \(C(x)\) is
locally constant (piecewise constant with jumps only at the threshold
crossings, which occur with probability zero). Hence \(C\) is
Lipschitz-almost-everywhere (trivially, with Lipschitz constant 0 on
each continuity region). \(\square\)

\textbf{Implications for the proof chain:} - The Lyapunov descent (Lemma
SE-1.1) uses \(\mathbb{E}[\cdot]\), which is unchanged by null-set
modifications. - The SCX update
\(\text{SCXUpdate}(S_t, \mathcal{M}_t, f_{\theta_t})\) inherits
Lipschitz-almost-everywhere continuity because it aggregates \(C(x)\)
over finite memory banks. - The stochastic approximation theory applies
to the mean field, which is Lipschitz continuous. - Therefore:
\textbf{the indicator discontinuity does NOT break the proof chain}. The
gap identified in the ARXIV verification is closed by this lemma.

\textbf{Proposition SE-1.0 (Lipschitz Composition).} Under C2 and C3,
the student loss function is Lipschitz in \(\theta\) with composite
constant \(L_\ell \cdot L_f\), where \(L_\ell\) is the Lipschitz
constant of the loss \(\ell\):

\[
|\ell(f_{\theta_1}(x), y) - \ell(f_{\theta_2}(x), y)| \leq L_\ell L_f \|\theta_1 - \theta_2\|
\]

and the gradient is bounded by \(G = L_\ell L_f\).

\subsubsection{3.3 Conditions C4-C6: Learning Rates and
Sampling}\label{conditions-c4-c6-learning-rates-and-sampling}

Condition C4 ensures enough exploration (\(\sum \alpha_t = \infty\))
while controlling noise (\(\sum \alpha_t^2 < \infty\)). Condition C5
ensures that the memory bank \(\mathcal{M}_t\) is a representative
sample of the acceptance-biased distribution. Condition C6 ensures that
the system eventually stabilizes rather than oscillating indefinitely
due to non-vanishing step sizes.

\subsubsection{3.4 Condition C7: Bounded Gatekeeper
Update}\label{condition-c7-bounded-gatekeeper-update}

This ensures that the gatekeeper scores remain in \([0,1]\) and that the
update increments are integrable. The projection \(\Pi_{[0,1]}\)
enforces this explicitly.

\begin{center}\rule{0.5\linewidth}{0.5pt}\end{center}

\subsection{4. Conclusion of Theorem
SE-1}\label{conclusion-of-theorem-se-1}

Under conditions (C1)-(C7), the sequence \((S_t, \theta_t)\) converges
to a \textbf{fixed point} \((S^*, \theta^*)\) almost surely, where:

\textbf{(a) Self-consistent gatekeeper.} \(S^*\) is self-consistent in
the Bayesian sense:

\[
S^*(x,y) = P(\text{clean} \mid x, y, \mathcal{M}_\infty), \quad \forall (x,y) \in \mathcal{X} \times \mathcal{Y}
\]

where \(\mathcal{M}_\infty = \lim_{t \to \infty} \mathcal{M}_t\) is the
limiting memory bank (possibly infinite). Equivalently, \(S^*\)
satisfies the fixed-point equation:

\[
S^* = \text{SCXUpdate}(S^*, \mathcal{M}_\infty, f_{\theta^*})
\]

\textbf{(b) Optimal NEP student (local).} \(\theta^*\) is a local
minimum of the expected NEP loss under the limiting data distribution
induced by \(S^*\):

\[
\theta^* = \arg\min_{\theta \in \Theta} \mathbb{E}_{(x,y) \sim P_{S^*}}[\ell(f_\theta(x), y)] \quad \text{(local)}
\]

where \(P_{S^*}(x,y) \propto S^*(x,y) \cdot P_0(x,y)\) is the
acceptance-biased distribution at the fixed point.

\textbf{(c) Self-consistency equation.} The fixed point satisfies the
coupled system:

\[
\begin{aligned}
S^* &= \mathcal{T}_S(S^*, \theta^*) \\
\theta^* &= \mathcal{T}_\theta(S^*, \theta^*)
\end{aligned}
\]

where \(\mathcal{T}_S\) and \(\mathcal{T}_\theta\) are the gatekeeper
and student update operators respectively, evaluated at the fixed point.

\textbf{(d) Lyapunov convergence.} The total loss function
\(V(S, \theta) = \mathbb{E}_{(x,y) \sim P_S}[\ell(f_\theta(x), y)]\)
serves as a Lyapunov function that is strictly decreasing along the
trajectory until the fixed point is reached.

\begin{center}\rule{0.5\linewidth}{0.5pt}\end{center}

\subsection{5. Lemma SE-1.1: Lyapunov Function
Non-Increase}\label{lemma-se-1.1-lyapunov-function-non-increase}

\textbf{Lemma SE-1.1 (Lyapunov Function).} Define the total system loss:

\[
V_t = V(S_t, \theta_t) = \mathbb{E}_{(x,y) \sim P_{S_t}}[\ell(f_{\theta_t}(x), y)] + \lambda \cdot \|S_t - S^*_{\text{prior}}\|^2_{\mathcal{L}^2}
\]

where \(\lambda > 0\) regularizes the gatekeeper toward a prior
\(S^*_{\text{prior}}\) (e.g., the initial SCX score), and \(P_S\) is the
acceptance-biased distribution under gatekeeper \(S\).

Under conditions (C2)-(C7) and \textbf{C10} (reference-set replay ---
REQUIRED for Lyapunov descent; without C10 Lemma SE-1.1 is NOT
guaranteed, see Theorem 12.2):

\[
\mathbb{E}[V_{t+1} \mid \mathcal{F}_t] \leq V_t - \eta_t
\]

where \(\eta_t \geq 0\) is a non-negative random variable satisfying:

\[
\eta_t = \alpha_t \cdot \|\nabla_{\theta} L_t(\theta_t)\|^2 + \beta_t \cdot \|\Delta S_t\|^2_{\mathcal{L}^2} + o(\alpha_t + \beta_t)
\]

with
\(\Delta S_t = \text{SCXUpdate}(S_t, \mathcal{M}_t, f_{\theta_t}) - S_t\).

\emph{Proof sketch.}

\begin{enumerate}
\def\labelenumi{\arabic{enumi}.}
\item
  \textbf{Decompose the change.} The total change in \(V_t\) has three
  components:

  \begin{itemize}
  \tightlist
  \item
    Change due to \(\theta\) update:
    \(\Delta_\theta V_t = V(S_t, \theta_{t+1}) - V(S_t, \theta_t)\)
  \item
    Change due to \(S\) update:
    \(\Delta_S V_t = V(S_{t+1}, \theta_{t+1}) - V(S_t, \theta_{t+1})\)
  \item
    Cross terms: \(o(\alpha_t + \beta_t)\)
  \end{itemize}
\item
  \textbf{Student contribution.} From the Lipschitz gradient assumption
  (SE-A7):

  \[
  \mathbb{E}[\Delta_\theta V_t \mid \mathcal{F}_t] \leq -\alpha_t \|\nabla_\theta L_t(\theta_t)\|^2 + \frac{L_g \alpha_t^2}{2} G^2
  \]
\item
  \textbf{Gatekeeper contribution.} From the Lipschitz property of
  SCXUpdate:

  \[
  \mathbb{E}[\Delta_S V_t \mid \mathcal{F}_t] \leq -\beta_t \left( \|\Delta S_t\|^2_{\mathcal{L}^2} - 2\lambda \langle \Delta S_t, S_t - S^*_{\text{prior}} \rangle \right) + O(\beta_t^2)
  \]

  By the Cauchy-Schwarz inequality and the bound
  \(\|\Delta S_t\|_\infty \leq B_S\), the cross term is \(O(\beta_t)\).
\item
  \textbf{Combining terms.} Under sufficient annealing
  (\(\beta_t \to 0\), \(\alpha_t \to 0\)), the negative quadratic terms
  dominate the \(O(\alpha_t^2 + \beta_t^2)\) corrections, giving the
  supermartingale inequality.
\end{enumerate}

\textbf{Corollary SE-1.1 (Convergence of \(V_t\)).} The Lyapunov
sequence \(\{V_t\}\) converges almost surely to a finite limit
\(V_\infty\):

\[
V_t \xrightarrow{a.s.} V_\infty \geq 0
\]

and moreover:

\[
\sum_{t=1}^\infty \bigl( \alpha_t \|\nabla_\theta L_t(\theta_t)\|^2 + \beta_t \|\Delta S_t\|^2_{\mathcal{L}^2} \bigr) < \infty \quad \text{a.s.}
\]

\subsubsection{5.1 Selection Bias Cycle Analysis (DEFECT-13
Fix)}\label{selection-bias-cycle-analysis-defect-13-fix}

A critical concern for the Lyapunov descent proof is the
\textbf{selection bias cycle} inherent in the self-evolution loop:

\[\text{Gatekeeper } S_t \xrightarrow{\text{selects}} \text{Memory } \mathcal{M}_{t+1} \xrightarrow{\text{trains}} \text{NEP } f_{\theta_{t+1}} \xrightarrow{\text{feedback}} \text{Gatekeeper } S_{t+1}\]

\textbf{The problem.} At each iteration, the gatekeeper \(S_t\) selects
which samples to admit to \(\mathcal{M}_{t+1}\) based on its current
scoring function. The NEP student is then trained on this
acceptance-biased dataset. The NEP's outputs feed back to update the
gatekeeper. This creates a self-reinforcing cycle:

\begin{enumerate}
\def\labelenumi{\arabic{enumi}.}
\tightlist
\item
  If \(S_t\) is overly confident about certain regions of
  \(\mathcal{X}\), it will admit samples predominantly from those
  regions, excluding counter-examples.
\item
  The NEP student, trained on this biased subset, will perform well on
  the admitted regions but may degrade on excluded regions.
\item
  The updated gatekeeper \(S_{t+1}\) receives feedback from a student
  that is biased toward the accepted distribution, reinforcing the
  original bias.
\end{enumerate}

\textbf{Formalization.} Let \(P_0\) be the original data distribution
and \(P_{S_t}\) be the acceptance-biased distribution:

\[P_{S_t}(x,y) \propto \mathbf{1}\{S_t(x,y) \geq \gamma\} \cdot P_0(x,y).\]

The Lyapunov function
\(V_t = \mathbb{E}_{(x,y) \sim P_{S_t}}[\ell(f_{\theta_t}(x), y)]\)
measures performance on the \textbf{acceptance-biased} distribution, not
on the original distribution \(P_0\). A decrease in \(V_t\) does
\textbf{not} guarantee improvement on \(P_0\) --- it may simply reflect
the gatekeeper becoming better at selecting easy samples.

\textbf{Quantifying the distribution shift.} The total variation between
the original and acceptance-biased distributions at time \(t\) is:

\[TV(P_0, P_{S_t}) = \frac{1}{2} \sum_{x,y} |P_0(x,y) - P_{S_t}(x,y)| \leq \mathbb{P}_{P_0}(S_t(x,y) < \gamma).\]

As the gatekeeper becomes more selective (higher \(\gamma_t\) or
stricter \(S_t\)), the distribution shift can grow, making the Lyapunov
decrease on \(P_{S_t}\) increasingly disconnected from true performance
on \(P_0\).

\textbf{Mitigation strategies (proposed for Conjecture SE-1):}

\begin{enumerate}
\def\labelenumi{\arabic{enumi}.}
\item
  \textbf{Reference-set evaluation (DEFECT-03/04 fix).} Evaluate both
  gatekeeper and student on a fixed reference set \(M_0\) (held out
  before self-evolution begins), not on the growing memory bank:
  \[\Phi_{\text{ref}}(S_t, \theta_t) = \frac{1}{|M_0|} \sum_{x \in M_0} (S_t(x,y(x)) - \hat{C}(x))^2 + \frac{\lambda}{|V_0|} \sum_{(x,y) \in V_0} \ell(f_{\theta_t}(x), y).\]
\item
  \textbf{Random exploration (C9).} A fixed fraction \(\varepsilon > 0\)
  of accepted samples are chosen uniformly at random regardless of
  \(S_t\) score, ensuring all regions remain represented in
  \(\mathcal{M}_t\).
\item
  \textbf{Annealing acceptance threshold (C8).} The threshold
  \(\gamma_t\) starts low (permissive, \(\gamma_0 < 0.5\)) and increases
  to \(0.5\) as \(t \to \infty\):
  \[\gamma_t = \gamma_0 + (0.5 - \gamma_0)(1 - e^{-t/\tau_\gamma}).\]
  This prevents early over-confidence while allowing increasing
  selectivity as the gatekeeper improves.
\item
  \textbf{Two-timescale separation (DEFECT-14 fix).} Gatekeeper updates
  are asymptotically slower than student updates
  (\(\beta_t = o(\alpha_t)\)), ensuring the student converges to a local
  minimum on the current distribution before the distribution shifts
  significantly.
\end{enumerate}

\textbf{Degradation under bias accumulation.} If the selection bias is
not controlled, the effective information available to the system
degrades. Let
\(\mathcal{R}_{\text{excluded}}(t) = \{x : S_t(x,y) < \gamma_t\}\) be
the set of excluded samples. The student's generalization error on the
full distribution \(P_0\) satisfies:

\[\mathbb{E}_{P_0}[\ell(f_{\theta_t}(x), y)] = \mathbb{E}_{P_{S_t}}[\ell(f_{\theta_t}(x), y)] \cdot \mathbb{P}(S_t \geq \gamma_t) + \mathbb{E}_{\mathcal{R}_{\text{excluded}}}[\ell(f_{\theta_t}(x), y)] \cdot \mathbb{P}(S_t < \gamma_t).\]

The second term is unobserved if no exploration is performed, and can
grow without detection.

\textbf{Status.} The selection bias cycle is the primary obstacle to
proving Lyapunov descent. The mitigation strategies (C8, C9) are
\textbf{conjectured} to be sufficient, but a rigorous proof that they
break the self-reinforcing cycle remains open. Without these
mitigations, the Lyapunov decrease on \(P_{S_t}\) is \textbf{not} a
valid indicator of system improvement --- it may reflect distribution
shift, not genuine model improvement.

\begin{center}\rule{0.5\linewidth}{0.5pt}\end{center}

\subsection{6. Lemma SE-1.2: Finite Cover of Memory Bank
Types}\label{lemma-se-1.2-finite-cover-of-memory-bank-types}

\textbf{Lemma SE-1.2 (Finite Memory Bank Configurations).} Under
condition C1 (finite \(\mathcal{X}\)) and the acceptance criterion
\(S_t(x,y) \geq \gamma\):

\textbf{(a)} The number of possible memory bank configurations is
finite.

\textbf{(b)} The memory bank \(\mathcal{M}_t\) evolves monotonically
(\(\mathcal{M}_t \subseteq \mathcal{M}_{t+1}\)) and stabilizes after
finitely many additions:

\[
\mathcal{M}_t = \mathcal{M}_{t+1} = \cdots = \mathcal{M}_\infty \quad \text{for all sufficiently large } t
\]

\textbf{(c)} Consequently, a.s. there exists a finite time \(T\) such
that for all \(t \geq T\):

\[
\Delta \mathcal{M}_t = \mathcal{M}_{t+1} \setminus \mathcal{M}_t = \varnothing
\]

\emph{Proof sketch.}

\begin{enumerate}
\def\labelenumi{\arabic{enumi}.}
\item
  \textbf{Finiteness.} Since \(\mathcal{X}\) is finite,
  \(\mathcal{X} \times \mathcal{Y}\) is finite. The memory bank
  \(\mathcal{M}_t\) is a subset of \(\mathcal{X} \times \mathcal{Y}\).
  There are at most \(2^{|\mathcal{X} \times \mathcal{Y}|}\) possible
  memory banks, which is finite.
\item
  \textbf{Monotonicity.} By definition,
  \(\mathcal{M}_{t+1} = \mathcal{M}_t \cup \text{accepted}_t\), so
  \(\mathcal{M}_t\) is monotonically increasing
  (\(\mathcal{M}_t \subseteq \mathcal{M}_{t+1}\)).
\item
  \textbf{Stabilization.} A monotonically increasing sequence in a
  finite set must stabilize after finitely many steps. Therefore there
  exists \(T\) such that \(\mathcal{M}_t = \mathcal{M}_{t+1}\) for all
  \(t \geq T\).
\end{enumerate}

\textbf{Corollary SE-1.2 (Infinite \(\mathcal{X}\) Extension).} If
\(\mathcal{X}\) is infinite but has finite covering dimension, the
number of ``types'' (equivalence classes of configurations under the SCX
score topology) is finite. The stabilization result holds up to
\(\varepsilon\)-equivalence: for any \(\varepsilon > 0\), there exists
\(T_\varepsilon\) such that:

\[
\sup_{x,y} |S_t(x,y) - S_{t'}(x,y)| < \varepsilon \quad \text{for all } t, t' \geq T_\varepsilon
\]

This extension requires the additional assumption that the SCX score
function is uniformly continuous in the data distribution.

\begin{center}\rule{0.5\linewidth}{0.5pt}\end{center}

\subsection{7. Lemma SE-1.3: Parameter Displacement
Vanishes}\label{lemma-se-1.3-parameter-displacement-vanishes}

\textbf{Lemma SE-1.3 (Vanishing Parameter Displacement).} Under
conditions (C2)-(C7):

\[
\|\theta_{t+1} - \theta_t\| \xrightarrow{a.s.} 0 \quad \text{and} \quad \|S_{t+1} - S_t\|_\infty \xrightarrow{a.s.} 0
\]

as \(t \to \infty\).

\emph{Proof sketch.}

\begin{enumerate}
\def\labelenumi{\arabic{enumi}.}
\item
  \textbf{Student displacement.} From the update rule and bounded
  gradient (SE-A8):

  \[
  \|\theta_{t+1} - \theta_t\| = \alpha_t \|g_t(\theta_t)\| \leq \alpha_t G
  \]

  Since \(\alpha_t \to 0\) (by C6),
  \(\|\theta_{t+1} - \theta_t\| \to 0\).
\item
  \textbf{Gatekeeper displacement.} From the update rule and bounded
  update (C7):

  \[
  \|S_{t+1} - S_t\|_\infty = \beta_t \|\text{SCXUpdate}(S_t, \mathcal{M}_t, f_{\theta_t}) - S_t\|_\infty \leq \beta_t B_S
  \]

  Since \(\beta_t \to 0\) (by C6), \(\|S_{t+1} - S_t\|_\infty \to 0\).
\end{enumerate}

\textbf{Corollary SE-1.3 (Joint Convergence).} The sequence
\((S_t, \theta_t)\) is Cauchy in the product space
\(\mathcal{L}^\infty(\mathcal{X} \times \mathcal{Y}) \times \mathbb{R}^{d_\theta}\)
with probability 1:

\[
\lim_{t \to \infty} \sup_{t' \geq t} \left( \|S_{t'} - S_t\|_\infty + \|\theta_{t'} - \theta_t\| \right) = 0 \quad \text{a.s.}
\]

\begin{center}\rule{0.5\linewidth}{0.5pt}\end{center}

\subsection{8. Lemma SE-1.4: Limit Points Are Fixed
Points}\label{lemma-se-1.4-limit-points-are-fixed-points}

\textbf{Lemma SE-1.4 (Fixed Point Characterization).} Let
\((S_\infty, \theta_\infty)\) be any limit point of the sequence
\((S_t, \theta_t)\) (along a subsequence \(t_k \to \infty\)). Then:

\textbf{(a) Gatekeeper fixed point.} \(S_\infty\) satisfies:

\[
S_\infty = \text{SCXUpdate}(S_\infty, \mathcal{M}_\infty, f_{\theta_\infty})
\]

\textbf{(b) Student fixed point.} \(\theta_\infty\) satisfies:

\[
\nabla_\theta \mathbb{E}_{(x,y) \sim P_{S_\infty}}[\ell(f_{\theta_\infty}(x), y)] = 0
\]

\textbf{(c) Cross-consistency.} The pair \((S_\infty, \theta_\infty)\)
is a joint fixed point of the coupled system:

\[
\begin{aligned}
S_\infty &= \mathcal{T}_S(S_\infty, \theta_\infty) \\
\theta_\infty &= \mathcal{T}_\theta(S_\infty, \theta_\infty)
\end{aligned}
\]

\emph{Proof sketch.}

\begin{enumerate}
\def\labelenumi{\arabic{enumi}.}
\item
  \textbf{Gatekeeper fixed point.} From Lemma SE-1.1 and Corollary
  SE-1.1:

  \[
  \sum_{t=1}^\infty \beta_t \|\Delta S_t\|^2_{\mathcal{L}^2} < \infty \quad \text{a.s.}
  \]

  Since \(\sum \beta_t = \infty\) (C4), it must be that
  \(\|\Delta S_t\|_{\mathcal{L}^2} \to 0\) along a subsequence with
  \(\beta_t > 0\). By the Lipschitz continuity of SCXUpdate (C3):

  \[
  \|S_\infty - \text{SCXUpdate}(S_\infty, \mathcal{M}_\infty, f_{\theta_\infty})\|_\infty = 0
  \]
\item
  \textbf{Student fixed point.} From Lemma SE-1.1 and Corollary SE-1.1:

  \[
  \sum_{t=1}^\infty \alpha_t \|\nabla_\theta L_t(\theta_t)\|^2 < \infty \quad \text{a.s.}
  \]

  Since \(\sum \alpha_t = \infty\) (C4),
  \(\|\nabla_\theta L_t(\theta_t)\| \to 0\) along a subsequence. By
  Lemma SE-1.3 and the continuity of \(\nabla L_\infty\), convergence of
  \(\theta_t\) implies \(\nabla_\theta L_\infty(\theta_\infty) = 0\).
\item
  \textbf{Cross-consistency.} From Lemma SE-1.2, \(\mathcal{M}_t\)
  stabilizes at \(\mathcal{M}_\infty\). At the limit point, \(S_\infty\)
  induces \(P_{S_\infty}\), \(\theta_\infty\) minimizes the loss under
  \(P_{S_\infty}\), and \(S_\infty\) is the Bayes estimate under
  \(\mathcal{M}_\infty\). These three conditions are mutually
  consistent.
\end{enumerate}

\begin{center}\rule{0.5\linewidth}{0.5pt}\end{center}

\subsection{9. Proof Sketch of Theorem
SE-1}\label{proof-sketch-of-theorem-se-1}

Here we assemble the lemmas into the complete proof of Theorem SE-1.

\textbf{Step 1: Lyapunov descent (Lemma SE-1.1).} Show that
\(V_t = V(S_t, \theta_t)\) is a supermartingale with non-negative drift.
From the update rules, the one-step expected change in \(V_t\) is
bounded above by \(-\eta_t\) where \(\eta_t \geq 0\) measures progress.
By the supermartingale convergence theorem, \(V_t \to V_\infty\) a.s.
and \(\sum \eta_t < \infty\) a.s.

\textbf{Step 2: Memory bank stabilization (Lemma SE-1.2).} Because
\(\mathcal{X}\) is finite (C1), the memory bank \(\mathcal{M}_t\) can
take only finitely many values. Since it is monotonically increasing
(\(\mathcal{M}_t \subseteq \mathcal{M}_{t+1}\)), it must stabilize after
finitely many steps. Thus there exists \(T < \infty\) such that
\(\mathcal{M}_{t} = \mathcal{M}_\infty\) for all \(t \geq T\).

\textbf{Step 3: Displacement vanishes (Lemma SE-1.3).} From the
annealing condition (C6), \(\alpha_t \to 0\) and \(\beta_t \to 0\), so
successive iterates are arbitrarily close:

\[
\|\theta_{t+1} - \theta_t\| \leq \alpha_t G \to 0, \quad \|S_{t+1} - S_t\|_\infty \leq \beta_t B_S \to 0
\]

\textbf{Step 4: Convergence of parameter sequence.} The Lyapunov descent
(Step 1) combined with vanishing displacement (Step 3) implies that the
sequence \((S_t, \theta_t)\) is Cauchy. More precisely:

\begin{itemize}
\tightlist
\item
  From Step 1: \(\sum \alpha_t \|\nabla L_t(\theta_t)\|^2 < \infty\) and
  \(\sum \beta_t \|\Delta S_t\|^2 < \infty\).
\item
  From Step 3: the displacements vanish.
\item
  For any \(\varepsilon > 0\), choose \(t\) large enough that
  \(\max(\alpha_t, \beta_t) < \varepsilon\), then the total remaining
  displacement from \(t\) onward is bounded by
  \(\sum_{s \geq t} \alpha_s G + \beta_s B_S\), which can be made
  arbitrarily small by further increasing \(t\) (since
  \(\alpha_s, \beta_s \to 0\) and the convergence of \(V_t\) implies the
  optimization progress is finite).
\end{itemize}

Thus \((S_t, \theta_t)\) converges to a limit
\((S_\infty, \theta_\infty)\) almost surely.

\textbf{Step 5: Limit is a fixed point (Lemma SE-1.4).} Having
established convergence, we verify that the limit satisfies the fixed
point equations:

\[
\begin{aligned}
S_\infty &= \text{SCXUpdate}(S_\infty, \mathcal{M}_\infty, f_{\theta_\infty}) \\
\nabla_\theta \mathbb{E}_{P_{S_\infty}}[\ell(f_{\theta_\infty})] &= 0
\end{aligned}
\]

These follow from: - By Step 2, \(\mathcal{M}_t = \mathcal{M}_\infty\)
for all large \(t\), so the SCX update is applied to a fixed memory
bank. By the Lyapunov descent and
\(\sum \beta_t \|\Delta S_t\|^2 < \infty\), the SCX update must converge
to the identity. - By the Robbins-Monro convergence (Theorem 5.1)
applied to the fixed distribution \(P_{S_\infty}\), \(\theta_t\)
converges to a stationary point of the expected loss.

\textbf{Step 6: Structure of the fixed point.} The fixed point satisfies
the self-consistency conditions stated in the theorem:

\begin{itemize}
\tightlist
\item
  \(S_\infty(x,y) = P(\text{clean} \mid x, y, \mathcal{M}_\infty)\)
  because the SCX update with infinite data converges to the true
  conditional probability (by the Glivenko-Cantelli theorem and the
  consistency of the SCX estimator).
\item
  \(\theta_\infty\) is a local minimum because the gradient vanishes and
  the Hessian is positive semi-definite at any stationary point of a
  smooth loss (assuming the loss has no saddle points or that the
  algorithm escapes them).
\end{itemize}

\textbf{Remark on the proof.} Step 2 relies crucially on the finite
structure space (C1). For infinite \(\mathcal{X}\), the memory bank may
never stabilize, leading to regimes 2-4 below. The full proof for
infinite \(\mathcal{X}\) requires additional compactness arguments and
is partially conjectured (see Section 14).

\begin{center}\rule{0.5\linewidth}{0.5pt}\end{center}

\subsection{10. Four Convergence
Regimes}\label{four-convergence-regimes}

While Theorem SE-1 establishes convergence under conditions (C1)-(C7),
violations of these conditions lead to distinct long-term behaviors. We
characterize four regimes:

\subsubsection{Regime 1: Classical Convergence (Conditions C1-C7 fully
satisfied)}\label{regime-1-classical-convergence-conditions-c1-c7-fully-satisfied}

\begin{itemize}
\tightlist
\item
  \((S_t, \theta_t) \to (S^*, \theta^*)\) almost surely
\item
  \(\mathcal{M}_t \to \mathcal{M}_\infty\) finite
\item
  Monotonic improvement: \(V_{t+1} \leq V_t\) in expectation
\item
  The system reaches a self-consistent configuration
\item
  \textbf{Occurs when}: Structure space is finite, gatekeeper
  stabilizes, student converges to a local minimum
\end{itemize}

\textbf{Formal guarantee.} Under conditions (C1)-(C7), Theorem SE-1
provides:

\[
\mathbb{P}\left( \lim_{t \to \infty} (S_t, \theta_t) = (S^*, \theta^*) \right) = 1
\]

\subsubsection{Regime 2: Limit Cycle (C4/C6 partially
violated)}\label{regime-2-limit-cycle-c4c6-partially-violated}

\begin{itemize}
\tightlist
\item
  \((S_t, \theta_t)\) does not converge to a single point but oscillates
  among a finite set of configurations
  \(\{(S^{(1)}, \theta^{(1)}), \dots, (S^{(K)}, \theta^{(K)})\}\)
\item
  The memory bank \(\mathcal{M}_t\) oscillates correspondingly
\item
  The Lyapunov function \(V_t\) does not converge but cycles within a
  bounded range
\item
  \textbf{Occurs when}: \(\alpha_t\) or \(\beta_t\) do not vanish
  (constant learning rates), the loss landscape has multiple deep
  basins, and the noise in gradient estimates causes periodic jumps
  between basins
\end{itemize}

\textbf{Proposition SE-1.1 (Limit Cycle Conditions).} If
\(\alpha_t = \alpha\) (constant) and \(\beta_t = \beta\) (constant), and
the product \(\alpha \beta\) exceeds a threshold, the system enters a
limit cycle with period \(T\) satisfying:

\[
S_{t+T} = S_t, \quad \theta_{t+T} = \theta_t, \quad \forall t \text{ sufficiently large}
\]

\emph{Proof sketch.} The system becomes a deterministic discrete-time
dynamical system with noise. For sufficiently large step sizes, the
update
\((\theta, S) \mapsto (\theta - \alpha \nabla L_S(\theta), S + \beta \Delta(S, \theta))\)
may have periodic orbits. By the finiteness of \(\mathcal{X}\) (C1), the
state space of \((\theta, S)\) can be partitioned into finitely many
cells, and by the pigeonhole principle, any infinite bounded trajectory
must repeat a state.

\textbf{Conjectured extension.} For infinite \(\mathcal{X}\), limit
cycles may be replaced by more complex attractors (strange attractors,
quasi-periodic orbits). This is an open theoretical question.

\subsubsection{Regime 3: Perpetual Discovery (C1 violated in data
distribution)}\label{regime-3-perpetual-discovery-c1-violated-in-data-distribution}

\begin{itemize}
\tightlist
\item
  \(\mathcal{M}_t\) grows without bound, constantly incorporating new
  types of data
\item
  \(S_t\) never stabilizes because the data distribution is
  non-stationary
\item
  \(\theta_t\) tracks a moving target
\item
  \textbf{Occurs when}: The data distribution \(P_0\) has infinite
  support \textgreater{} the memory bank capacity, and the gatekeeper
  continues to find novel reliable samples
\end{itemize}

\textbf{Proposition SE-1.2 (Perpetual Discovery Bound).} Under perpetual
discovery, for any \(T < \infty\):

\[
\liminf_{t \to \infty} \frac{1}{t} \sum_{i=1}^t \|\mathcal{M}_{i+1} \setminus \mathcal{M}_i\| > 0
\]

i.e., new samples are added at a non-vanishing rate. In this regime:

\begin{itemize}
\tightlist
\item
  The gatekeeper does not converge in \(\ell^\infty\) norm:
  \(\limsup \|S_{t+1} - S_t\|_\infty > 0\)
\item
  The student's generalization error is bounded below by the irreducible
  error from distribution shift (Theorem 5.6)
\item
  The total variation of the data distribution \(P_t\) is unbounded:
  \(\sum_{t=1}^\infty \text{TV}(P_{t+1}, P_t) = \infty\)
\end{itemize}

\textbf{Conjectured bound.} Under perpetual discovery, the expected
student loss satisfies:

\[
\liminf_{t \to \infty} \mathbb{E}[\ell(f_{\theta_t}(x), y)] \geq \inf_{\theta} \mathbb{E}_{P_\infty}[\ell(f_\theta(x), y)] + \tau \cdot \limsup \text{TV}(P_t, P_\infty)
\]

where \(\tau\) is a constant depending on the loss Lipschitz constant
and the student capacity. This represents the irreducible error from
never converging to a fixed distribution.

\subsubsection{Regime 4: Divergent Collapse (S\_t becomes
degenerate)}\label{regime-4-divergent-collapse-s_t-becomes-degenerate}

\begin{itemize}
\tightlist
\item
  \(S_t(x,y) \to 0\) for most/all \((x,y)\), causing the gatekeeper to
  reject all data
\item
  \(\mathcal{M}_t\) stops growing (or shrinks if a removal mechanism
  exists)
\item
  The student receives no training signal and diverges
\item
  \textbf{Occurs when}: The SCX update systematically underestimates
  reliability, the acceptance threshold \(\gamma\) is set too high, or
  noise overwhelms the system
\end{itemize}

\textbf{Proposition SE-1.3 (Divergent Collapse Condition).} Divergent
collapse occurs if there exists a time \(T\) such that:

\[
S_t(x,y) < \gamma, \quad \forall (x,y) \in \mathcal{X} \times \mathcal{Y}, \forall t \geq T
\]

i.e., the gatekeeper rejects all inputs. Under this condition:

\begin{itemize}
\tightlist
\item
  \(\mathcal{M}_t = \mathcal{M}_T\) for all \(t \geq T\) (no new data)
\item
  \(\theta_t\) evolves according to
  \(\theta_{t+1} = \theta_t - \alpha_t g_t(\theta_t)\) on a fixed finite
  dataset, which converges to a stationary point (by standard RM theory)
\item
  \(S_t\) updates without new data, converging to the fixed point of the
  update equation on \(\mathcal{M}_T\)
\end{itemize}

\textbf{Prevention mechanisms.} The following mechanisms prevent
divergent collapse:

\begin{enumerate}
\def\labelenumi{\arabic{enumi}.}
\tightlist
\item
  \textbf{Exploration bonus}: Add a small positive constant to \(S_t\):
  \(\tilde{S}_t = (1-\varepsilon) S_t + \varepsilon\).
\item
  \textbf{Threshold annealing}: Reduce \(\gamma_t \to 0\) as
  \(t \to \infty\).
\item
  \textbf{Optimistic initialization}: Initialize
  \(S_0(x,y) = 1 - \delta\) for small \(\delta\).
\item
  \textbf{Memory replay}: Periodically inject random samples regardless
  of \(S_t\) score.
\end{enumerate}

\textbf{Empirical observation.} Divergent collapse is rarely observed in
practice because the SCX update uses a Bayesian averaging that ensures
\(S_t\) remains bounded away from zero for at least some regions of the
input space.

\begin{center}\rule{0.5\linewidth}{0.5pt}\end{center}

\subsection{11. Conditions Separating the
Regimes}\label{conditions-separating-the-regimes}

The following table summarizes the conditions that determine which
regime occurs:

\begin{longtable}[]{@{}
  >{\raggedright\arraybackslash}p{(\linewidth - 10\tabcolsep) * \real{0.0941}}
  >{\raggedright\arraybackslash}p{(\linewidth - 10\tabcolsep) * \real{0.1765}}
  >{\raggedright\arraybackslash}p{(\linewidth - 10\tabcolsep) * \real{0.1412}}
  >{\raggedright\arraybackslash}p{(\linewidth - 10\tabcolsep) * \real{0.1294}}
  >{\raggedright\arraybackslash}p{(\linewidth - 10\tabcolsep) * \real{0.3059}}
  >{\raggedright\arraybackslash}p{(\linewidth - 10\tabcolsep) * \real{0.1529}}@{}}
\toprule\noalign{}
\begin{minipage}[b]{\linewidth}\raggedright
Regime
\end{minipage} & \begin{minipage}[b]{\linewidth}\raggedright
\(\mathcal{X}\)
\end{minipage} & \begin{minipage}[b]{\linewidth}\raggedright
\(\alpha_t\)
\end{minipage} & \begin{minipage}[b]{\linewidth}\raggedright
\(\beta_t\)
\end{minipage} & \begin{minipage}[b]{\linewidth}\raggedright
\(\mathcal{M}_t\) behavior
\end{minipage} & \begin{minipage}[b]{\linewidth}\raggedright
\(S_t\) limit
\end{minipage} \\
\midrule\noalign{}
\endhead
\bottomrule\noalign{}
\endlastfoot
1. Classical & Finite & \(\to 0\), RM & \(\to 0\) & Stabilizes & Fixed
point \(S^*\) \\
2. Limit cycle & Finite & Constant & Constant & Periodic & Periodic
orbit \\
3. Perpetual discovery & Infinite & \(\to 0\) or constant & \(\to 0\) or
constant & Grows unbounded & Does not converge \\
4. Divergent collapse & Any & Any & Any & Stops growing &
\(S_t \approx 0\) everywhere \\
\end{longtable}

\textbf{Phase diagram (conjectured).} The behavior is governed by two
key parameters:

\begin{enumerate}
\def\labelenumi{\arabic{enumi}.}
\tightlist
\item
  \textbf{Exploration rate}:
  \(\rho_{\text{explore}} = \lim_{t \to \infty} \frac{|\mathcal{M}_t \setminus \mathcal{M}_{t-1}|}{|\mathcal{M}_{t-1}|}\)
  (fraction of new samples added per step)
\item
  \textbf{Learning rate ratio}:
  \(r = \lim_{t \to \infty} \alpha_t / \beta_t\) (relative speed of
  student vs.~gatekeeper)
\end{enumerate}

The conjectured phase diagram is:

\[
\text{Regime} = \begin{cases}
1 \text{ (Classical)} & \text{if } \rho_{\text{explore}} = 0 \text{ and } r \to 0 \\
2 \text{ (Limit cycle)} & \text{if } \rho_{\text{explore}} = 0 \text{ and } r \text{ is bounded away from } 0 \\
3 \text{ (Discovery)} & \text{if } \rho_{\text{explore}} > 0 \\
4 \text{ (Collapse)} & \text{if } S_t \text{ threshold never met}
\end{cases}
\]

This phase diagram is \textbf{conjectured} --- a rigorous analysis of
the phase transitions is an open problem.

\begin{center}\rule{0.5\linewidth}{0.5pt}\end{center}

\subsection{12. Connection to Theorem 1 (Noise
Detection)}\label{connection-to-theorem-1-noise-detection}

\textbf{Proposition SE-1.4 (Noise Detection During Self-Evolution).}
Throughout the self-evolution process (all four regimes), the noise
detection guarantee of Theorem 1 applies at each time step \(t\):

\[
\text{F1}_t \geq 1 - \frac{1}{\eta} \sum_{s \in \mathcal{S}} \rho_s \cdot \exp\!\bigl(-2M\Delta_s(t)^2\bigr)
\]

where \(\Delta_s(t)\) is the state-\(s\) separation gap evaluated using
the current estimated state probabilities \(\rho_s\) and expert risks
\(R_m(s)\).

\emph{Proof.} Theorem 1's assumptions (A1)-(A6) are about the
\textbf{data generation process} and \textbf{expert training}, not about
the self-evolution loop. As long as these assumptions hold for the base
data and experts, the noise detection guarantee holds at every time
step. The self-evolution loop may change the effective noise
distribution (via memory bank filtering), but Theorem 1's bound applies
to the \textbf{original} noise detection task, which is invariant.

\textbf{Corollary SE-1.5 (Monotonic Noise Detection Improvement).} If
the self-evolution successfully improves the state estimates
\(\rho_s^{(t)}\) and expert risk estimates \(R_m^{(t)}(s)\), the
effective \(\Delta_s(t)\) may increase over time, leading to tighter
noise detection:

\[
\Delta_s(t) \geq \Delta_s(0) \quad \text{for sufficiently large } t
\]

This is \textbf{not guaranteed} --- if the gatekeeper biases the sample
distribution away from certain states, the state estimates may degrade.
In Regime 4 (collapse), \(\Delta_s(t) \to 0\).

\textbf{Implication.} The noise detection guarantee provides a
\textbf{lower bound} on performance that holds regardless of
self-evolution. The self-evolution can only improve (or leave unchanged)
this bound; it cannot make the bound worse than the initial guarantee.

\textbf{Remark on Correct Global F1 Aggregation (Lemma F Fix).}
Proposition SE-1.4 uses Theorem 1's bound, which correctly operates on
the \textbf{global confusion matrix}:

\[TP_{\text{global}} = \eta N \sum_s \rho_s \cdot TPR_s,\quad
FP_{\text{global}} = (1-\eta)N \sum_s \rho_s \cdot FPR_s,\quad
FN_{\text{global}} = \eta N \sum_s \rho_s \cdot FNR_s.\]

The global F1 score is
\(F1_{\text{global}} = 2 \cdot TP_{\text{global}} / (2 \cdot TP_{\text{global}} + FP_{\text{global}} + FN_{\text{global}})\).
The Hoeffding bound applies directly to \(FPR_{\text{global}}\) and
\(FNR_{\text{global}}\) because these error rates are \textbf{linear} in
per-state quantities
(\(FPR_{\text{global}} = \sum_s \rho_s \cdot FPR_s\),
\(FNR_{\text{global}} = \sum_s \rho_s \cdot FNR_s\)). The F1 algebraic
inequality (Lemma B) then yields the stated bound without requiring
additivity of F1 itself.

\begin{quote}
\textbf{Corrected from original Lemma F (DEFECT-01/02).} The original
claim \(F1_{\text{global}} = \sum_s \rho_s \cdot F1_s\) is
\textbf{mathematically false}: F1 is a harmonic mean, not a linear
expectation. The correct relationship expands \(1 - F1_{\text{global}}\)
as
\(\frac{1}{2}FNR_{\text{global}} + \frac{1-\eta}{2\eta}FPR_{\text{global}} + o(e^{-M\kappa_{\min}})\),
where the linearity of FNR and FPR allows state-wise summation. States
with \(\kappa_s > \kappa_{\min}\) contribute
\(o(e^{-M(\kappa_s-\kappa_{\min})})\) and are asymptotically negligible.
\end{quote}

\textbf{Remark on Bahadur-Rao Lattice Correction (DEFECT-06).} For
Bernoulli-distributed expert errors (lattice span \(h=1\)), the correct
Bahadur-Rao expansion uses the factor \((1-e^{-\lambda^*})^{-1}\) rather
than \(1/\lambda^*\):

\[\mathbb{P}(\bar{X}_M \geq \theta) = \frac{\exp(-M \cdot KL(\theta\|p))}{(1-e^{-\lambda^*}) \sqrt{2\pi M \cdot \theta(1-\theta)}} \bigl(1 + O(1/M)\bigr),\]

where \(\lambda^* = \log\frac{\theta(1-p)}{p(1-\theta)}\). The
correction factor \(\lambda^*/(1-e^{-\lambda^*})\) ranges from
\(\sim 1.3\) to \(3.2\) for typical \(\lambda^* \in [0.5, 3]\). The
corrected minimax constant \(C_{\min}^{\text{(corr)}}\) differs from the
original \(C_{\min}\) by a multiplicative factor
\(\mathcal{L}(\lambda_0^*, \lambda_1^*)\) of approximately
\(1.05\)--\(1.12\) for typical parameter ranges. \textbf{The minimax
optimality conclusion is unchanged} because both the achievability bound
and the lower bound receive the same lattice correction factor --- the
constant matching \(C_{\text{SCX}}/C_{\min}^{\text{(corr)}} = 1 + o(1)\)
is preserved. Only the \textbf{claimed numerical value} of the constants
changes.

\textbf{Corrected \(C_{\min}\) formula:}
\[C_{\min}^{\text{(corr)}} = \frac{\eta}{2} \left(\frac{1-\eta}{\eta}\right)^{s} \frac{1/\lambda_0^* + 1/|\lambda_1^*|}{\sqrt{\theta^*(1-\theta^*)}} \times \mathcal{L}(\lambda_0^*, \lambda_1^*),\]
where the lattice correction factor is:
\[\mathcal{L}(\lambda_0^*, \lambda_1^*) = \frac{\lambda_0^* |\lambda_1^*|}{(1-e^{-\lambda_0^*})(1-e^{-|\lambda_1^*|})} \cdot \frac{\frac{1-e^{-\lambda_0^*}}{\lambda_0^*} + \frac{1-e^{-|\lambda_1^*|}}{|\lambda_1^*|}}{\frac{1}{\lambda_0^*} + \frac{1}{|\lambda_1^*|}}.\]

For Test Case 1 parameters (\(p_0=0.10\), \(p_1=0.60\)):
\(\lambda_0^* = 1.404\), \(|\lambda_1^*| = 1.198\),
\(\mathcal{L} \approx 1.085\) (an \(\sim 8.5\%\) change).

\begin{center}\rule{0.5\linewidth}{0.5pt}\end{center}

\subsection{13. Connection to Theorem 3
(Unidentifiability)}\label{connection-to-theorem-3-unidentifiability}

\textbf{Proposition SE-1.5 (Corrected --- Non-Resolution of
Unidentifiability).} Theorem 3 establishes a fundamental
population-level unidentifiability between noise and intrinsic
difficulty: there exist two data-generating processes
\(\mathcal{P}_{\text{noise}}\) and \(\mathcal{P}_{\text{hard}}\) with
identical observable joint distributions, such that any algorithm has
expected error at least \(\eta\rho/2\) on the ambiguous subset,
\textbf{regardless of sample size}.

The self-evolution loop \textbf{does not resolve this
unidentifiability}, because:

\begin{enumerate}
\def\labelenumi{\arabic{enumi}.}
\tightlist
\item
  The two worlds produce identical distributions at \textbf{all} sample
  sizes:
  \(TV(\mathcal{P}_{\text{noise}}^n, \mathcal{P}_{\text{hard}}^n) = 0\)
  for every \(n \geq 1\). By the i.i.d. product property, if the
  single-sample distributions are identical, so are all finite-sample
  product distributions.
\item
  The \(\eta\rho/2\) lower bound is a \textbf{population-level minimax
  bound} that holds for \(n = \infty\), not a finite-sample bound that
  decays.
\item
  No amount of data accumulation through the memory bank can distinguish
  distributions that are identical at every finite sample size.
\end{enumerate}

\textbf{Why the original \(1/\sqrt{N_t}\) rate was incorrect.} The
original Proposition SE-1.5 confused two distinct problems: -
\textbf{Testing two simple hypotheses with TV \textgreater{} 0}: Error
decays as \(O(1/\sqrt{N})\) (standard Chernoff-Stein regime). -
\textbf{Testing two hypotheses with TV = 0} (Theorem 3's construction):
Error is \textbf{constant} with \(N\) --- the two hypotheses produce
exactly the same distribution, so no test can do better than random
guessing on the ambiguous subset.

Since Theorem 3's construction has
\(TV(\mathcal{P}_{\text{noise}}, \mathcal{P}_{\text{hard}}) = 0\), the
error is \(\eta\rho/2\) for all \(N\), including \(N \to \infty\). The
claimed \(1/\sqrt{N_t}\) decay does not apply.

\textbf{What self-evolution CAN improve.} However, the self-evolution
loop can \textbf{improve state estimation} (via Theorem 5's cluster
consistency) and \textbf{reduce expert error rates} (via better feature
representations from the NEP student). These improvements indirectly
tighten the \textbf{bounding constants} in Theorem 1's F1 guarantee.
They are \textbf{orthogonal} to the unidentifiability in Theorem 3: they
improve performance on the \emph{identifiable} portion of the data but
cannot resolve the \emph{unidentifiable} portion.

\textbf{Corollary SE-1.6 (Corrected).} The asymptotic bound
\(\lim_{t\to\infty} \text{Error}_{\text{minimax}}(t) = 0\) claimed in
the original version does \textbf{not} hold. The correct statement is:

\[\liminf_{t\to\infty} \text{Error}_{\text{minimax}}(t) \geq \frac{\eta\rho}{2} > 0,\]

i.e., the irreducible error from unidentifiability persists even with
infinite self-evolution data.

\emph{Proof.} Theorem 3's construction has
\(TV(\mathcal{P}_{\text{noise}}, \mathcal{P}_{\text{hard}}) = 0\),
implying
\(TV(\mathcal{P}_{\text{noise}}^n, \mathcal{P}_{\text{hard}}^n) = 0\)
for all \(n\) by the i.i.d. product property. The lower bound
\(\eta\rho/2\) is derived from a minimax argument over the two worlds
and does \textbf{not} depend on \(n\). The self-evolution loop cannot
create information that is absent from the single-sample distribution.
\(\square\)

\begin{center}\rule{0.5\linewidth}{0.5pt}\end{center}

\subsection{14. What is Proven
vs.~Conjectured}\label{what-is-proven-vs.-conjectured}

\begin{longtable}[]{@{}
  >{\raggedright\arraybackslash}p{(\linewidth - 4\tabcolsep) * \real{0.3182}}
  >{\raggedright\arraybackslash}p{(\linewidth - 4\tabcolsep) * \real{0.3636}}
  >{\raggedright\arraybackslash}p{(\linewidth - 4\tabcolsep) * \real{0.3182}}@{}}
\toprule\noalign{}
\begin{minipage}[b]{\linewidth}\raggedright
Claim
\end{minipage} & \begin{minipage}[b]{\linewidth}\raggedright
Status
\end{minipage} & \begin{minipage}[b]{\linewidth}\raggedright
Notes
\end{minipage} \\
\midrule\noalign{}
\endhead
\bottomrule\noalign{}
\endlastfoot
Theorem SE-1: Convergence to fixed point & \textbf{Proven under C1-C7} &
Relies on finite structure space (C1) \\
Theorem SE-1: Self-consistency equation & \textbf{Proven} & Follows from
fixed point condition \\
Lemma SE-1.1: Lyapunov non-increase & \textbf{CONDITIONAL} (requires Thm
12.5 reference-set replay) & Supermartingale argument valid only with
replay mechanism; without replay, Thm 12.2 proves Lyapunov descent is
formally impossible \\
Lemma SE-1.2: Finite memory bank types & \textbf{Proven} & Finite
\(\mathcal{X}\) guarantee \\
Lemma SE-1.3: Parameter displacement vanishes & \textbf{Proven} &
Annealing condition \\
Lemma SE-1.4: Limit points are fixed points & \textbf{Proven} & From
Lemmas SE-1.1 through SE-1.3 \\
Regime 1: Classical convergence & \textbf{Proven} & Theorem SE-1 \\
Regime 2: Limit cycle & \textbf{Proven for finite \(\mathcal{X}\),
constant step sizes} & Periodic orbits in finite-state DDS \\
Regime 3: Perpetual discovery & \textbf{Conjectured} & Infinite
\(\mathcal{X}\) extension, rate bound is heuristic \\
Regime 4: Divergent collapse & \textbf{Proven conditionally} & If
\(S_t < \gamma\) for all \((x,y)\), the dynamics are determined \\
Proposition SE-1.1: Limit cycle conditions & \textbf{Proven} & Finite
state space + constant step sizes \\
Proposition SE-1.2: Perpetual discovery bound & \textbf{Conjectured} &
The irreducibility bound lacks rigorous proof \\
Proposition SE-1.3: Collapse condition & \textbf{Proven} & Direct
consequence of update rules \\
Phase diagram & \textbf{Conjectured} & Phase transitions are not
rigorously characterized \\
Proposition SE-1.4: Noise detection during evolution & \textbf{Proven} &
Theorem 1 applies invariantly \\
Proposition SE-1.5: Resolution of unidentifiability & \textbf{RETRACTED}
(was ``Proven'') & DEFECT-05: The \(1/\sqrt{N_t}\) rate was based on a
confusion between TV \textgreater{} 0 and TV = 0 cases. Theorem 3's
construction has TV = 0 at all \(N\), so the error is constant
(\(\eta\rho/2\)), not decaying. Corrected version now acknowledges
non-resolution. \\
Infinite \(\mathcal{X}\) extension of Lemma SE-1.2 &
\textbf{Conjectured} & Requires uniform continuity and compactness \\
\end{longtable}

\textbf{Summary of gaps.} The main open theoretical questions are:

\begin{enumerate}
\def\labelenumi{\arabic{enumi}.}
\item
  \textbf{Infinite \(\mathcal{X}\) convergence}: Whether the memory bank
  stabilizes (or only stabilizes up to \(\varepsilon\)) for infinite
  input spaces depends on the topology of \(\mathcal{X}\) and the
  continuity of \(S_t\).
\item
  \textbf{Perpetual discovery rate}: The precise relationship between
  the data distribution's tail behavior and the discovery rate is not
  fully characterized.
\item
  \textbf{Phase transitions}: The boundaries between regimes 1-3 and the
  nature of the transitions (sharp vs.~smooth, first-order
  vs.~second-order) are open questions.
\item
  \textbf{Finite-sample bounds}: Theorem SE-1 is an asymptotic result.
  Finite-sample convergence rates are not established.
\end{enumerate}

\begin{center}\rule{0.5\linewidth}{0.5pt}\end{center}

\subsection{15. References}\label{references}

\begin{enumerate}
\def\labelenumi{\arabic{enumi}.}
\tightlist
\item
  Robbins, H., \& Monro, S. (1951). A stochastic approximation method.
  \emph{Annals of Mathematical Statistics}, 22(3), 400-407.
\item
  Doob, J. L. (1953). \emph{Stochastic Processes}. Wiley.
\item
  Kushner, H. J., \& Yin, G. G. (2003). \emph{Stochastic Approximation
  and Recursive Algorithms and Applications} (2nd ed.). Springer.
\item
  Borkar, V. S. (2008). \emph{Stochastic Approximation: A Dynamical
  Systems Viewpoint}. Cambridge University Press.
\item
  Polyak, B. T. (1987). \emph{Introduction to Optimization}.
  Optimization Software.
\item
  Bertsekas, D. P., \& Tsitsiklis, J. N. (1996). \emph{Neuro-Dynamic
  Programming}. Athena Scientific.
\item
  Khalil, H. K. (2002). \emph{Nonlinear Systems} (3rd ed.). Prentice
  Hall.
\item
  Strogatz, S. H. (2018). \emph{Nonlinear Dynamics and Chaos} (2nd ed.).
  CRC Press.
\item
  Ghosh, J. K., \& Ramamoorthi, R. V. (2003). \emph{Bayesian
  Nonparametrics}. Springer.
\item
  van der Vaart, A. W. (1998). \emph{Asymptotic Statistics}. Cambridge
  University Press.
\item
  Kleijn, B. J. K., \& van der Vaart, A. W. (2012). The Bernstein-von
  Mises theorem under misspecification. \emph{Electronic Journal of
  Statistics}, 6, 354-381.
\item
  Castillo, I., \& Nickl, R. (2013). Nonparametric Bernstein-von Mises
  theorems in Gaussian white noise. \emph{Annals of Statistics}, 41(4),
  1999-2028.
\end{enumerate}

\begin{center}\rule{0.5\linewidth}{0.5pt}\end{center}

\emph{End of 06\_fixed\_point\_convergence.md -- Central theorem of SCX
self-evolution convergence.}

\begin{center}\rule{0.5\linewidth}{0.5pt}\end{center}

\subsection{Changelog}\label{changelog}

\begin{longtable}[]{@{}
  >{\raggedright\arraybackslash}p{(\linewidth - 6\tabcolsep) * \real{0.1875}}
  >{\raggedright\arraybackslash}p{(\linewidth - 6\tabcolsep) * \real{0.2500}}
  >{\raggedright\arraybackslash}p{(\linewidth - 6\tabcolsep) * \real{0.2500}}
  >{\raggedright\arraybackslash}p{(\linewidth - 6\tabcolsep) * \real{0.3125}}@{}}
\toprule\noalign{}
\begin{minipage}[b]{\linewidth}\raggedright
Date
\end{minipage} & \begin{minipage}[b]{\linewidth}\raggedright
Defect
\end{minipage} & \begin{minipage}[b]{\linewidth}\raggedright
Change
\end{minipage} & \begin{minipage}[b]{\linewidth}\raggedright
Severity
\end{minipage} \\
\midrule\noalign{}
\endhead
\bottomrule\noalign{}
\endlastfoot
2026-06-28 & DEFECT-01/02 & \textbf{Fixed Lemma F}: Added correct global
F1 aggregation derivation (Section 12 remark). The original claim
\(F1_{\text{global}} = \sum_s \rho_s \cdot F1_s\) was false due to F1's
nonlinearity. Corrected: FNR and FPR are linear in per-state quantities;
the F1 expansion absorbs the nonlinearity into
\(o(e^{-M\kappa_{\min}})\). Theorem 1's existing bound is already
correct. & FATAL \\
2026-06-28 & DEFECT-03/04 & \textbf{Fixed Lyapunov function}: Updated
cross-reference to the explicit reference-set-based \(\Phi\) defined in
Document 02 (Section 7). The Lyapunov convergence claim is downgraded to
Conjecture; descent property is conditional on two-timescale separation
(C6'), random exploration (C8), and annealing threshold (C9). & FATAL \\
2026-06-28 & DEFECT-05 & \textbf{Retracted SE-1.5}: Replaced the false
\(1/\sqrt{N_t}\) resolution rate claim with the correct analysis:
Theorem 3's unidentifiability is population-level (\(TV=0\) at all
\(N\)), so the \(\eta\rho/2\) error is constant, not decaying. Corollary
SE-1.6 corrected to \(\liminf \geq \eta\rho/2 > 0\). Self-evolution can
improve bounding constants but not resolve unidentifiability. & FATAL \\
2026-06-28 & DEFECT-06 & \textbf{Added Bahadur-Rao lattice correction}:
Bernoulli lattice correction \((1-e^{-\lambda^*})^{-1}\) replaces
\(1/\lambda^*\) in all asymptotic formulas. \(C_{\min}\) changes by
\(\sim 5\)--\(12\%\); minimax optimality preserved (both bounds get same
correction). & MAJOR \\
2026-06-28 & DEFECT-13 & \textbf{Added selection bias cycle analysis}
(Section 5.1): Gatekeeper selects data → NEP trains on biased data → NEP
output feeds back to gatekeeper. Lyapunov decrease on \(P_{S_t}\) may
reflect distribution shift, not genuine improvement. Added proposed
mitigations: reference-set evaluation, random exploration (C8),
annealing threshold (C9). & MAJOR \\
2026-06-28 & DEFECT-14 & \textbf{Fixed exploration-stabilization
tension} (C4/C6): Replaced original C6
(\(\max(\alpha_t,\beta_t) \to 0\)) with two-timescale condition
\(\beta_t = o(\alpha_t)\). Added variance control
\(\sum \beta_t^2 < \infty\) for gatekeeper updates. Provided canonical
scheduling scheme \(\alpha_t = t^{-0.6}\), \(\beta_t = t^{-0.8}\)
satisfying all joint conditions. & MAJOR \\
2026-06-28 & B1 & \textbf{Lemma SE-1.1 status changed to CONDITIONAL}:
Updated §14 table and C10 note to reflect that Lemma SE-1.1 (Lyapunov
non-increase) is conditional on reference-set replay (Thm 12.5). Without
replay, Lemma SE-1.1 is not guaranteed (Thm 12.2). Fixes
self-contradiction between Doc 06 and Doc 10. & FATAL \\
2026-06-28 & B2 & \textbf{Added C10 to condition list}: Reference-set
replay is now an explicit required condition (\(\mathcal{M}_0\),
\(\mathcal{V}_0\) held out before self-evolution). The Lyapunov function
\(\Psi\) is evaluated on fixed reference sets, breaking the selection
bias cycle. C10 bridges Lemma SE-1.1 with Theorems 12.2/12.5. & FATAL \\
2026-06-28 & B2 (SE-1 body) & \textbf{Added reference-set note to SE-1
statement body}: The gatekeeper update
\(\text{SCXUpdate}(S_t, \mathcal{M}_{t+1}, f_{\theta_{t+1}})\) is
evaluated on a fixed reference set \(\mathcal{M}_0\) (see C10). This
anchors evaluation to a stationary distribution and prevents trivial
loss reduction via sample selection. & MAJOR \\
\end{longtable}

\end{document}
